\section*{INTRODUÇÃO}

No monitoramento de incêndios florestais, o registro pontual de uma anomalia térmica apresenta utilidade limitada para a tomada de decisão quando desprovido de contexto espacial e temporal imediato. Para subsidiar ações de campo, a literatura postula a necessidade de construir \textit{Consciência Situacional} (ENDSLEY, 1995) --- a capacidade de perceber, compreender e projetar o efeito de eventos dinâmicos sobre comunidades e infraestruturas do território. Embora o país conte com dados de sensoriamento remoto (INPE, 2026), as informações sobre infraestruturas e vulnerabilidades permanecem distribuídas em bases setoriais descentralizadas (ONS, DATASUS, DNIT, FUNAI), dificultando cruzamentos em tempo hábil.

Diante desse cenário, este projeto avalia o desenvolvimento de uma plataforma analítica geoespacial orientada a eventos. A proposta concentra-se em correlacionar anomalias térmicas e dados de infraestrutura crítica em tempo quase real, buscando prover suporte informacional contínuo para a Defesa Civil e operadores setoriais.
